% Mathematical fragment for the raw clausal computational-equality report.
% Needs amsmath, amssymb, amsthm and theorem/lemma/corollary environments.
% No normalization operation, external solver, or function variables are used.

\section{A raw clausal calculus with five inference patterns}
\label{sec:ce5-specification}

Write $\mathsf{CE5}$ for the following presentation. Numerical terms are
raw finite expressions
\[
 t ::= x\mid 0\mid S(t)\mid I(n,a;u.h).
\]
Only $u$ in $h$ is bound. All substitution is ordinary capture-avoiding
substitution. In particular, the terms $I(0,a;u.h)$ and $a$ are distinct
syntax, including after substitution. Formulas have just two constructors:
\[
 F::=(s=t)\mid D_{x,y}F.
\]
The two coordinates of $D$ are bound. All displayed binders may be
renamed before a substitution.

A context is a finite reusable list of \emph{clause occurrences}
\[
 [A;y.B(y)],\qquad y\notin\operatorname{FV}(A).
\]
The clause binds $y$ only in its consequent. It is a context object,
not a third formula constructor. Put $H=(0=0)$. An ordinary formula
assumption $F$ abbreviates the constant clause
\[
 \operatorname{ord}(F):=[H;v.F],
 \qquad v\notin\operatorname{FV}(F).
\]
Thus $\Gamma,F$ below means $\Gamma,\operatorname{ord}(F)$.
Clause occurrences, rather than only their values, are retained: a
clause-use instruction selects an occurrence in the context.

\subsection{The partial computational equality operation}

For raw numerical terms $s,t,b$, define $\Phi(s,t,b)$ by these three
cases, and leave it undefined otherwise:
\[
\begin{array}{c|c|c}
 b&\hbox{required literal source}&\Phi(s,t,b)\\ \hline
 s&b\equiv s&t\\
 I(0,a;u.h)&s\equiv a&t\\
 I(Sn,a;u.h)&s\equiv I(n,a;u.h)&h[t/u].
\end{array}
\]
Here $\equiv$ is literal identity of well-scoped finite syntax, modulo
the chosen representation of bound names. The rows are disjoint. In the
second row the source is a proper subterm of $b$; in the third row $b$
has one more successor node than its required source. Hence neither row
can also satisfy $b\equiv s$. The last two rows have different counter
constructors.

This operation performs one specified equality/computation step. It
does not normalize terms. Its last row uses equality of the whole
preceding iterator value, not equality of an initializer.

\subsection{Every primitive inference}

The complete inventory consists of the following five patterns:
\[
 \frac{}{\Gamma\vdash s=s}.                              \tag{Refl}
\]
\[
 \frac{\Gamma\vdash s=t\qquad\Gamma\vdash b=r}
      {\Gamma\vdash r=\Phi(s,t,b)}
 \quad\hbox{if $\Phi(s,t,b)$ is defined}.                \tag{Eval}
\]
\[
 \frac{\Gamma\vdash A}{\Gamma\vdash B(t)}
 \quad\hbox{if the occurrence $[A;y.B(y)]$ is in $\Gamma$}.
                                                               \tag{U}
\]
\[
 \frac{\Gamma\vdash D_{x,y}B(x,y)\qquad
       \Gamma,[B(b,0);y.B(b,y)]\vdash F}
      {\Gamma\vdash F}.                                 \tag{E}
\]
\[
 \frac{\Gamma,B(t,\eta)\vdash B(t,Sz)}
      {\Gamma\vdash D_{x,y}B(x,y)},\qquad
 \eta=
 \begin{cases}
 z,&B\text{ has equation root and }t\equiv0,\\
 0,&\text{otherwise}.
 \end{cases}                                           \tag{I}
\]
In E, $b$ is fresh for the retained context, the conclusion, and the
passive parameters of the matrix. In I, $z$ is fresh for the retained
context, witness $t$, and passive matrix parameters. In every instance,
the supplied terms and formulas must be well scoped. There is no
restriction that $t$ be a numeral.

No further logical, equality, successor, computation, or induction axiom
is primitive. In particular, neither a separate hypothesis rule nor a
separate $Q$ constructor occurs. The designated contradiction is $0=S0$.
The clause representation and the three cases of $\Phi$ are part of the
trusted specification and must be included in a complete checker.

\begin{lemma}[Ordinary assumptions and elementary proof operations]
Ordinary assumption use is admissible. So are weakening, injective fresh
renaming, and cut for ordinary formula assumptions.
\end{lemma}
\begin{proof}
Refl proves $H$ in every context. U on the selected occurrence
$[H;v.F]$, with any numerical target, then proves $F$.
Weakening and fresh renaming follow by induction on the finite
derivation. Fresh renaming preserves both the equation-root test and
the literal-zero test in I.

For ordinary-assumption cut, suppose a distinguished context occurrence
is $[H;v.F]$ and a proof of $F$ is supplied without that occurrence.
Induct on the derivation that uses it. At a U node selecting the
distinguished occurrence, replace the node by the supplied proof,
weakened into the local context. At every other node translate its
premises recursively. Eigenvariables can first be freshened.
This argument uses no arbitrary numerical-substitution theorem.
\end{proof}

\section{Recovering equality, computation, and the logical interface}
\label{sec:ce5-recovery}

The elementary equality constructions below are inherited from the
computational-equality proposal~\cite{computational}; they are repeated to make the present
rule inventory independently checkable. They work under arbitrary
retained clauses.

\begin{lemma}[Equality algebra and both computation equations]
Refl and Eval derive symmetry, transitivity, and
\[
 I(0,a;u.h)=a,\qquad
 I(Sn,a;u.h)=h[I(n,a;u.h)/u].
\]
\end{lemma}
\begin{proof}
Let $R_s$ denote Refl at $s$, and let $\operatorname{ev}(p,q)$ denote
Eval on its two supplied proofs. Its root case is
\[
 \frac{s=t\qquad s=r}{r=t}.
\]
For $p:s=t$, $\operatorname{ev}(R_s,p)$ proves $t=s$.
For $p:s=t$ and $q:t=v$,
$\operatorname{ev}(q,\operatorname{ev}(R_s,p))$ proves $s=v$.

For zero computation put $A=I(0,a;u.h)$; then
$\operatorname{ev}(R_a,R_A)$ proves $A=a$.
For successor computation put $C=I(n,a;u.h)$ and
$B=I(Sn,a;u.h)$; then
$\operatorname{ev}(R_C,R_B)$ proves $B=h[C/u]$.
Each computation proof has two reflexivity leaves and one Eval node.
\end{proof}

\begin{lemma}[Unrestricted congruence from one iteration]
\label{lem:ce5-congruence}
For every fixed numerical term template $h(u)$, the rule
\[
 \frac{\Gamma\vdash s=t}
      {\Gamma\vdash h[s/u]=h[t/u]}
\]
is admissible. The template may contain any number of occurrences of
$u$, including occurrences inside nested iterator bodies.
\end{lemma}
\begin{proof}
For the supplied proof $p:s=t$, set
$A_s=I(0,s;u.h)$ and $B_s=I(S0,s;u.h)$. Construct
\[
\begin{aligned}
 z_s&=\operatorname{ev}(R_s,R_{A_s})
       &&: A_s=s,\\
 c_s&=\operatorname{ev}(z_s,R_{B_s})
       &&: B_s=h[s/u],\\
 z_{s,t}&=\operatorname{ev}(p,R_{A_s})
       &&: A_s=t,\\
 \operatorname{Cong}_h(p)&=\operatorname{ev}(z_{s,t},c_s)
       &&: h[s/u]=h[t/u].
\end{aligned}
\]
The first and third lines use the zero row of $\Phi$. The second and
fourth use its successor row, whose required source is literally
$A_s$. Thus the construction assumes no prior congruence principle.
The body substitution performed by Eval supplies all the displayed
occurrences of $s$ and $t$.

There are four new Refl leaves and four Eval nodes, and the input proof
$p$ occurs once. This is a bound on inference nodes, not on the size of
the copied term labels or the serialized certificate.
\end{proof}

\begin{corollary}[One-term and equation-template transport]
The following rules are admissible, with no restriction on the source
term:
\[
 \frac{s=t\qquad h(s)=r}{r=h(t)},\qquad
 \frac{s=t\qquad a(s)=b(s)}{a(t)=b(t)}.
\]
\end{corollary}
\begin{proof}
For the first rule, apply the root case of Eval to the congruence
equation from Lemma~\ref{lem:ce5-congruence} and the second premise.
For the second, use the first rule twice:
\[
 s=t,\ a(s)=b(s)\ \Longrightarrow\ b(s)=a(t)
 \ \Longrightarrow\ a(t)=b(t).
\]
The second use takes the term template $b$.
\end{proof}

\begin{lemma}[All case introductions]
\label{lem:ce5-cases}
For every matrix $B$ and every witness $t$, the rule
\[
 \frac{\Gamma,B(t,0)\vdash B(t,Sz)}
      {\Gamma\vdash D_{x,y}B(x,y)}
\]
is admissible, with the usual freshness of $z$.
\end{lemma}
\begin{proof}
The primitive I rule already has this form except when $B$ is an
equation and $t$ is literally $0$. In that remaining case put
$c=I(0,0;u.0)$. The term $c$ is not literally zero, and zero computation
proves $c=0$.

Under the local assumption $B(c,0)$, equation-template transport gives
$B(0,0)$. By ordinary-assumption cut insert the supplied branch and
obtain $B(0,Sz)$. Symmetry and equation-template transport give
$B(c,Sz)$. Primitive I with witness $c$ is in its case mode, and its
conclusion is the required $D_{x,y}B(x,y)$. The conclusion has not
changed; only the proof's witnessing term has changed.
\end{proof}

\begin{lemma}[Transport in every formula]
\label{lem:ce5-formula-transport}
For each finite Eq/$D$ formula template $F(v)$, from
$\Gamma\vdash s=t$ and $\Gamma\vdash F(s)$ one can derive
$\Gamma\vdash F(t)$.
\end{lemma}
\begin{proof}
Induct externally on the formula template. Equation templates were
handled above. Suppose
$F(v)=D_{x,y}B(v,x,y)$. Eliminate the supplied $F(s)$ at a fresh
parameter $b$, retaining the clause
\[
 [B(s,b,0);y.B(s,b,y)].
\]
Introduce $F(t)$ by Lemma~\ref{lem:ce5-cases} with witness $b$.
Its branch assumes $B(t,b,0)$. By the induction hypothesis and the
symmetry of $s=t$, derive $B(s,b,0)$. U at $Sz$ gives
$B(s,b,Sz)$; the forward induction hypothesis gives $B(t,b,Sz)$.
Close the introduction and elimination. Every recursive transport
concerns the strictly smaller matrix $B$, so the proof is not circular.
\end{proof}

\subsection{Conditional closure without a formula constructor}

For a fresh unused $w$, abbreviate
\[
 q_yB(y):=D_{w,y}B(y).
\]
The dummy coordinate is a numerical binder, not a new sort or a domain
restriction.

\begin{lemma}[The three closure interfaces]
\label{lem:ce5-q}
The following operations are derivable:
\begin{enumerate}
\item From $\Gamma,B(0)\vdash B(Sz)$, derive $\Gamma\vdash q_yB$.
\item From $\Gamma\vdash q_yB$ and $\Gamma\vdash B(0)$, derive
      $\Gamma\vdash B(t)$ for every term $t$.
\item In any context retaining $[B(0);y.B(y)]$, derive $q_yB$.
\end{enumerate}
\end{lemma}
\begin{proof}
For the first operation use case introduction with unused witness.
For the second eliminate $q_yB$; its fresh witness does not occur in
$B$, and the resulting clause is exactly $[B(0);y.B(y)]$. U with the
supplied base proves $B(t)$; close E.
For the third, under the local case-introduction assumption $B(0)$,
use the retained clause at $Sz$, then close the first operation.
\end{proof}

For an equation $E(y)$, primitive I at dummy witness $0$ derives
$q_yE$ from $\Gamma,E(z)\vdash E(Sz)$. With its base and the second
closure interface this is contextual induction for equations.
For an equation $E(x,y)$ and any fixed term $t$, the same construction
first proves $q_yE(t,y)$ and then, by case introduction at $t$, recovers
the older witnessed progress conclusion $D_{x,y}E(x,y)$.
Numerical cases for arbitrary formulas come from the first two closure
interfaces. These are derived rules, not additional kernel instructions.

\section{Exact translation to and from the raw six-rule calculus}
\label{sec:ce5-equivalence}

Let $\mathsf{E6}$ be the raw computational-equality calculus already
present in the repository~\cite{computational}. Its formulas are Eq/$Q$/$D$, its contexts
are ordinary retained formulas, and its six patterns are Hyp, Refl,
Eval, QE, DE, and the same witness-directed DI. Refl and Eval are
exactly the operations above. Hyp selects a formula assumption, and
\[
 \frac{\Gamma\vdash Q_yB\qquad\Gamma\vdash B(0)}
      {\Gamma\vdash B(t)}
 \quad\text{is QE},
\]
\[
 \frac{\Gamma\vdash D_{x,y}B\qquad
       \Gamma,Q_yB(b,y)\vdash F}{\Gamma\vdash F}
 \quad\text{is DE}.
\]
DI has the displayed I conclusion and branch, with ordinary formula
assumptions instead of their constant clauses. The same freshness
conditions apply.

Define the translation $\tau$ on native formulas by
\[
 \tau(s=t)=(s=t),\qquad
 \tau(D_{x,y}B)=D_{x,y}\tau(B),\qquad
 \tau(Q_yB)=D_{w,y}\tau(B),
\]
where each newly inserted $w$ is fresh and unused. Numerical terms are
unchanged. Let $j$ include Eq/$D$ formulas into Eq/$Q$/$D$ formulas.
The translation $\tau$ preserves free numerical variables, commutes
with capture-free numerical substitution up to bound renaming, and
preserves the distinction between equation-root and compound formulas.

\begin{theorem}[Proof preservation and common-language equivalence]
\label{thm:ce5-proof-equivalence}
For every finite native context $\Gamma$ and native formula $A$,
\[
 \Gamma\vdash_{\mathsf{E6}} A
 \quad\Longrightarrow\quad
 \operatorname{ord}(\tau\Gamma)\vdash_{\mathsf{CE5}}\tau(A).
\]
For all finite Eq/$D$ contexts $\Delta$ and Eq/$D$ formulas $F$,
\[
 \operatorname{ord}(\Delta)\vdash_{\mathsf{CE5}}F
 \quad\Longleftrightarrow\quad
 j\Delta\vdash_{\mathsf{E6}}jF.
\]
The transformations are effective transformations of finite proof trees
and preserve their numerical terms.
\end{theorem}
\begin{proof}
For the first statement, induct on the $\mathsf{E6}$ proof.
Hyp becomes Refl at $0$ followed by U on the appropriate constant
clause. Refl and Eval are copied unchanged. QE becomes the second
interface of Lemma~\ref{lem:ce5-q}. DI is copied with translated
matrices and constant local clauses: $\tau$ preserves the matrix's
equation/compound root and leaves its numerical witness unchanged,
so the chosen mode is exactly the original mode.

For DE, first translate its major premise. E then opens
\[
 [\,\tau(B)(b,0);y.\tau(B)(b,y)\,].
\]
The third interface of Lemma~\ref{lem:ce5-q} derives
$q_y\tau(B)(b,y)$ in that context. This is exactly the translated old
branch assumption $\tau(Q_yB(b,y))$. Weaken the inductively translated
branch to the current context and cut in this proof of its ordinary
assumption. Closing E gives the translated conclusion. All inserted
dummy coordinates and eigenvariables are freshened first.

For the reverse statement, annotate each retained clause occurrence
by its introduction. An initial ordinary clause or an I branch clause
$[H;v.G]$ is represented in $\mathsf{E6}$ by the ordinary assumption
$jG$. A clause introduced by E is represented by the native assumption
$Q_yjB(b,y)$. If equal-looking clause values have different origins,
their occurrences retain their respective annotations.

Induct on the $\mathsf{CE5}$ proof under this invariant. A U selecting
an ordinary clause translates to Hyp; its proof of $H$ can be omitted.
A U selecting an E clause translates to QE on its associated $Q$
assumption and the translated antecedent proof. Refl and Eval are
identical. E translates to native DE, and I translates to native DI
with the same witness, matrix root, and local ordinary assumption.
This proves the reverse implication. The forward implication on the
common language is the first statement, since $\tau j$ is the identity.
\end{proof}

\begin{theorem}[Both formula composites]
\label{thm:ce5-composites}
For every Eq/$D$ formula $F$, $\tau jF=F$ literally, up to immaterial
bound renaming. For every native Eq/$Q$/$D$ formula $A$,
\[
 A\ \dashv\vdash_{\mathsf{E6}}\ j\tau A.
\]
Both numerical translations are identities.
\end{theorem}
\begin{proof}
The first claim follows from the definition. For the other one, the
essential native equivalence is
\[
 Q_yB(y)\ \dashv\vdash\
 D_{w,y}B(y),\qquad w\notin\operatorname{FV}(B).
\]
From $Q_yB$, assume $B(0)$, specialize the closure at $Sz$, and use
derived case introduction at an unused witness. Conversely, native DE
on the right-hand formula opens precisely $Q_yB$; Hyp closes the branch.

This equivalence extends through all enclosing binders. More
explicitly, a uniform pointwise interderivability of matrices lifts
through $D$ by DE followed by case introduction: convert the target
zero assumption back to the source zero instance, specialize the
opened source closure at $Sz$, and convert forward to the target
successor instance. It lifts through $Q$ by the same reasoning with
an unused witness. These constructions use only the already
established equality and case interfaces. External structural induction
on $A$ now proves the result.
\end{proof}

These proofs deliberately avoid claiming that whole proofs can be
substituted by reusing every rule tag. Although raw $\Phi$ commutes
with capture-free numerical substitution, a substitution can change
a nonliteral witness into literal $0$ and change I's mode. Full
substitution admissibility for ordinary retained-formula sequents
transfers from the native calculus through the equivalence above.
Weakening, fresh renaming and the ordinary cuts
used in the proof were established directly.

For arbitrary externally supplied clauses, use the ordinary reading
$[A;y.B]^\circ=A^\circ\to\forall yB^\circ(y)$. Its native representation
is the corresponding derived implication/universal macro. Clause use
then translates to implication elimination followed by universal
elimination. The exact common-language theorem does not require
external clauses beyond ordinary formula assumptions; it covers every
such retained-context consequence, not only closed theorems.


